% =====================================================================
% Limitations & conclusion.
% Adapted from RESEARCH_PLAN.md §11 point 7 (limitations) and §12
% (characterization framing).
% =====================================================================
\section{Limitations}
\label{sec:limitations}

Three limitations bound the interpretation of every result in this paper.

\paragraph{Batch size.}
The current harness runs with batch size~1. Batched prefill with
per-sample block-diagonal masks is what FlashAttention varlen is
designed for, and lifting the constraint would materially cut the cost
of the quality and efficiency sweeps. As stated, however, the reported
efficiency numbers are for the single-request setting; server-side
throughput at higher batch sizes is out of scope.

\paragraph{Preprint scope.}
This preprint covers retrieval quality (C1) on Qwen2.5-7B-Instruct, including
self-distillation finetuning to close the mask quality gap. TTFT / realized
speedup (C3) and KV-cache reuse (C4) are measured on the same model. The
restart position layout, which enables position-independent KV blocks for
caching, degrades generation quality even after mixed-layout finetuning;
this is a known limitation (see below) and the cache throughput claim is
scoped to the sequential layout.

\paragraph{Realization gap.}
If the analytical/realized gap of Section~\ref{sec:results-efficiency-gap}
does not close under the Phase~3 diagnostics --- if it is a structural
property of small-block varlen kernels on current hardware --- the
compute headline of DDI is smaller than the arithmetic suggests, and
the co-headline caching result of Section~\ref{sec:results-caching}
carries proportionally more of the paper. We flag this in the results
sections rather than only here, so a reader cannot miss the constraint.

\paragraph{Restart layout quality.}
The restart position layout (Section~\ref{sec:method-positions}) resets
each document's RoPE positions to zero, making KV blocks position-independent
and therefore cacheable across requests. On Qwen2.5-7B-Instruct, this layout
degrades generation quality (EM) relative to the sequential layout: 0.07 vs
0.39 on hotpot\_qa after mixed-layout finetuning. The perplexity gap closes
(1.628 vs 1.635) but the generation gap persists --- greedy decoding under
restart positions compounds small token-level errors. Mixed-layout training
(50/50 coin-flip per record) was attempted but each record sees only one
layout per epoch, so the model never learns consistent outputs across
layouts for the same record. The cache throughput claim (1.72$\times$ QPS)
is therefore scoped to the sequential layout. Future work: 2-epoch mixed
training, restart-only finetuning, or position-agnostic architecture changes.

\section{Conclusion}
\label{sec:conclusion}

We present a characterization of block-diagonal attention over retrieved
documents in decoder-only language models. On retrieval benchmarks we show
that \texttt{ddi\_qf} matches full-causal quality within the
non-inferiority margin after self-distillation finetuning
(Section~\ref{sec:results-quality}), report the realized versus analytical
prefill speedup at matched shapes (Section~\ref{sec:results-efficiency}),
and measure per-document KV-cache reuse throughput
(Section~\ref{sec:results-caching}). The caching path turns the
structural sparsity of the mask into a reusability property of KV
blocks, yielding 1.72$\times$ throughput on the sequential layout. The
restart position layout, designed to make KV blocks position-independent,
degrades generation quality even after mixed-layout finetuning; this
limitation is reported as a negative finding with root cause analysis.

The value of the paper does not depend on any variant beating a
baseline; it depends on the parity, efficiency, and caching axes being
measured honestly and reported with variance.
